\section{三角関数の積分}\label{節:三角式の整理}

$\sin x$ と $\cos x$ の原始関数は分かっていても,\,$\sin^2x$ や $\sin(3x)\cos x$ のような2乗や積は,基本公式にそのまま当てはめられない.そこで,三角関数の恒等式を使って $\sin(kx)$ や $\cos(kx)$ の和に直し,\sectionref{節:一次式の積分}の一次式の積分に帰着させよう.

2乗には二倍角の公式
\[
 \cos(2x)=1-2\sin^2x=2\cos^2x-1
\]
を使う.これを $\sin^2x$,\,$\cos^2x$ について解き直して
\[
 \sin^2x=\frac{1-\cos(2x)}2,\qquad
 \cos^2x=\frac{1+\cos(2x)}2
\]
とする.左辺は2乗だが,右辺は定数と余弦の和になった.

積には加法定理を使う.\,$\sin(\alpha+\beta)$ と $\sin(\alpha-\beta)$ の式を並べて足すと,
\[
 \begin{array}{rrl}
  &\sin(\alpha+\beta)={}&\sin\alpha\cos\beta+\cos\alpha\sin\beta\\
  +\,)&\sin(\alpha-\beta)={}&\sin\alpha\cos\beta-\cos\alpha\sin\beta\\ \hline
  &\sin(\alpha+\beta)+\sin(\alpha-\beta)={}&2\sin\alpha\cos\beta
 \end{array}
\]
となり,右辺の第2項が打ち消し合う.よって
\[
 \sin\alpha\cos\beta=\frac12\{\sin(\alpha+\beta)+\sin(\alpha-\beta)\}
\]
を得る.同様に,\,$\cos(\alpha-\beta)$ と $\cos(\alpha+\beta)$ の式を足し引きすると
\[
 \cos\alpha\cos\beta=\frac12\{\cos(\alpha-\beta)+\cos(\alpha+\beta)\},
\]
\[
 \sin\alpha\sin\beta=\frac12\{\cos(\alpha-\beta)-\cos(\alpha+\beta)\}
\]
となる.\,$\alpha=mx$,\,$\beta=nx$ とすれば,積が $\sin$ や $\cos$ の一次式の和に変わる.あとは各項を,角の $x$ の係数で割って積分すればよい.

\begin{bookpoint}
\textbf{POINT}\par
2乗は二倍角の公式で,積は加法定理を2つ足し引きして,\,$\sin(kx)$,\,$\cos(kx)$ の和に直し,\,$k$ で割って積分する.
\end{bookpoint}

\begin{example}{二乗と積を和へ直す}
次の積分を求めよ.
\[
 \text{(1)}\quad\int_0^{\frac\pi2}\sin^2x\dd x,
 \qquad
 \text{(2)}\quad\int\sin(3x)\cos x\dd x.
\]
\end{example}

\paragraphbf{方針}
\begin{enumerate}[label=(\arabic*)]
 \item 二倍角の公式で $\sin^2x=\frac{1-\cos(2x)}2$ と直す.
 \item 積を和に直す式で $\alpha=3x$,\,$\beta=x$ とする.
\end{enumerate}

\paragraphbf{解答}
\begin{enumerate}[label=(\arabic*)]
 \item $\frac\pi4$
 \item $-\frac18\cos(4x)-\frac14\cos(2x)+C$
\end{enumerate}

\paragraphbf{解法}
\begin{enumerate}[label=(\arabic*)]
 \item 二倍角の公式から
 \begin{align*}
  \int_0^{\frac\pi2}\sin^2x\dd x
  &=\frac12\int_0^{\frac\pi2}\{1-\cos(2x)\}\dd x\\
  &=\left[\frac x2-\frac14\sin(2x)\right]_0^{\frac\pi2}\\
  &=\frac\pi4
 \end{align*}
 となる.\,$\cos(2x)$ の積分で $2$ で割るので,全体の係数が $\frac14$ になる.定積分なので,任意定数は上下端で差を取ると消える.
 \item $\sin\alpha\cos\beta=\frac12\{\sin(\alpha+\beta)+\sin(\alpha-\beta)\}$ で $\alpha=3x$,\,$\beta=x$ とすると
 \[
  \sin(3x)\cos x=\frac12\{\sin(4x)+\sin(2x)\}
 \]
 となる.よって
 \[
  \int\sin(3x)\cos x\dd x
  =-\frac18\cos(4x)-\frac14\cos(2x)+C
 \]
 を得る.実数全体で成り立つ.右辺を微分すると $\frac12\sin(4x)+\frac12\sin(2x)$ となり,上の恒等式で積に戻る.
\end{enumerate}

\paragraphbf{公式を忘れたとき}
積を和へ直す3つの式を覚えていなくても,本節の初めのように加法定理を2つ書いて足し引きすれば,必要な式をその場で作れる.正弦と余弦の積なら正弦の加法定理を,正弦どうし,余弦どうしの積なら余弦の加法定理を並べればよい.

\paragraphbf{式の整理の限界}
本章では,積分記号の中の式を同じ値の別の式に書き換えてきた.どの変形を使うかは,元の式に現れた手がかりで決まる.整式の小さな積なら展開し,分子の次数が高ければ割り算をし,一次因子の積が分母にあれば部分分数を考える.二次式は平方完成でまとまりを作り,三角関数の2乗や積は恒等式で和へ直す.

ただし,式を長く展開すれば必ずよいわけではない.たとえば $x(x^2+1)^{20}$ は展開できるが,\,$x^2+1$ を1つの変数 $t$ と置けば計算を短くできる.そう置いてよい根拠を与えるのが,次章で学ぶ置換積分である.
